\documentclass[12pt, a4paper]{article}

% 引入必要的宏包
\usepackage[UTF8]{ctex}     % 中文支持
\usepackage{amsmath, amssymb, amsthm} % 数学符号和环境
\usepackage{geometry}       % 页面设置
\usepackage{enumitem}       % 列表定制
\usepackage{fancyhdr}       % 页眉页脚
\usepackage{cases}          % 分段函数增强

% 页面边距设置
\geometry{left=2.5cm, right=2.5cm, top=2.5cm, bottom=2.5cm}

% 宏定义：方便排版选择题选项
% 横向排列的选项 (4列)
\newcommand{\choicesFour}[4]{
    \begin{enumerate}[label=\Alph*., itemsep=0pt, parsep=0pt, topsep=5pt, labelsep=0.5em]
        \begin{minipage}{0.22\linewidth} \item #1 \end{minipage}
        \begin{minipage}{0.22\linewidth} \item #2 \end{minipage}
        \begin{minipage}{0.22\linewidth} \item #3 \end{minipage}
        \begin{minipage}{0.22\linewidth} \item #4 \end{minipage}
    \end{enumerate}
}
% 横向排列的选项 (2列)
\newcommand{\choicesTwo}[4]{
    \begin{enumerate}[label=\Alph*., itemsep=0pt, parsep=0pt, topsep=5pt, labelsep=0.5em]
        \begin{minipage}{0.45\linewidth} \item #1 \end{minipage}
        \begin{minipage}{0.45\linewidth} \item #2 \end{minipage}
        \begin{minipage}{0.45\linewidth} \item #3 \end{minipage}
        \begin{minipage}{0.45\linewidth} \item #4 \end{minipage}
    \end{enumerate}
}
% 纵向排列的选项 (1列)
\newcommand{\choices}[4]{
    \begin{enumerate}[label=\Alph*., itemsep=0pt, parsep=0pt, topsep=5pt, labelsep=0.5em]
        \item #1
        \item #2
        \item #3
        \item #4
    \end{enumerate}
}

% 填空题下划线
\newcommand{\blank}[1]{\underline{\hspace{#1}}}

\title{\textbf{《高等数学》必刷习题——提高版}}
\author{}
\date{}

\begin{document}

\maketitle

\section*{第一章 函数（提高）}

\begin{enumerate}
    \item 设 $ f(x)=\sin x $ , $ g(x)=\begin{cases}x-\pi, & x\leq0,\\x+\pi, & x>0,\end{cases} $ 则 $ f[g(x)]= $ \blank{1cm}。
    \choicesFour{$ \sin x $}{$ \cos x $}{$ - \sin x $}{$ - \cos x $}

    \item 函数 $ y = \ln | \sin x | $ 的定义域是 \blank{1cm}，其中 $k$ 为整数。
    \choicesTwo
    {$ x \neq \frac{k\pi}{2} $}
    {$ x \in (-\infty, \infty), x \neq k\pi $}
    {$ x = k\pi $}
    {$ x \in (-\infty, \infty) $}

    \item 函数 $ f(x)=\frac{1-x^{2}}{\sqrt{|x|-1}}+\arcsin(x-1) $ 的定义域是 \blank{1cm}。
    \choicesFour{$ (0,2] $}{$ [0,2] $}{$ (1,2] $}{$ [1,2] $}

    \item 函数 $ y=\sqrt{2x-x^{2}}-\arcsin\frac{2x-1}{7} $ 的定义域为 \blank{1cm}。
    \choicesFour{$ \left[-3,4\right] $}{$ (-3,4) $}{$ [0,2] $}{$ (0,2) $}

    \item 函数 $ y = \arcsin(1-x) + \frac{1}{2}\lg\frac{1+x}{1-x} $ 的定义域是 \blank{1cm}。
    \choicesFour{$ (0,1) $}{$ [0,1) $}{$ [0,1] $}{$ [0,1] $}

    \item 已知函数 $ f(x)=x^{3}+3x-2 $ , $ g(x)=\tan x $ , 则 $ f[g(\frac{\pi}{4})]= $ \blank{2cm}。

    \item 已知 $ f(x)=\begin{cases}\frac{1}{x}, & |x|>1\\0, & |x|\leq1\end{cases} $ 则 $ f[f(2021)]= $ \blank{2cm}。

    \item 函数 $ y = \ln[\ln(\ln x)] $ 的定义域是 \blank{2cm}。

    \item 函数 $ f(x)=\frac{1}{\sqrt{x-3}} $ 的定义域是 \blank{2cm}。

    \item 函数 $ f(x) = \sqrt{4 - x^{2}} + \frac{1}{\ln \cos x} $ 的定义域为 \blank{2cm}。

    \item 设 $ f(x-1)=x(x-1) $ ，则 $ f(x)= $ \blank{2cm}。

    \item 设 $ f(x) $ 是 $ (-\infty, +\infty) $ 内的偶函数，并且当 $ x \in (-\infty, 0) $ 时，有 $ f(x) = x + 2 $ ，则当 $ x \in (0, +\infty) $ 时， $ f(x) $ 的表达式是 \blank{2cm}。

    \item 函数 $ y = \sin x - \sin |x| $ 的值域是 \blank{2cm}。

    \item 设函数 $ f(x)=x^{3}-x^{2}-1 $ , 则 $ f[f(1)]= $ \blank{2cm}。

    \item 已知函数 $ f(x)=\frac{x+1}{x-1},x\in(1,+\infty) $ ，求复合函数 $ f[f(x)] $ 。
\end{enumerate}

\section*{第二章 极限（提高）}

\begin{enumerate}
    \item $ f(x) $ 在点 $ x = x_{0} $ 处有定义，是 $ f(x) $ 在 $ x = x_{0} $ 处连续的 \blank{1cm}。
    \choicesFour{必要条件}{充分条件}{充分必要条件}{无关条件}

    \item $ \lim_{x\to0}\frac{e^{x^{2}}-1}{\cos x-1}= $ \blank{1cm}。
    \choicesFour{$ \infty $}{2}{0}{-2}

    \item 下列极限正确的是 \blank{1cm}。
    \choicesTwo
    {$ \lim_{x\to0}\frac{1}{x}\sin x=1 $}
    {$ \lim_{x\to\infty}\frac{1}{x}\sin x=1 $}
    {$ \lim_{x\to0}x\sin\frac{1}{x}=1 $}
    {$ \lim_{x\to\infty}x\sin\frac{1}{x}=0 $}

    \item $ \lim_{x\to\infty}\frac{\cos e^{x}}{x}= $ \blank{1cm}。
    \choicesFour{$ \frac{\pi}{2} $}{0}{1}{$ \infty $}

    \item 以下极限中存在的是 \blank{1cm}。
    \choicesFour{$ \lim_{x\to0}2^{\frac{1}{x}} $}{$ \lim_{x\to0}\arctan\frac{1}{x} $}{$ \lim_{x\to0}x\sin\frac{1}{x} $}{$ \lim_{x\to0}\sin\frac{1}{x} $}

    \item 求极限 $ \lim_{x \to \infty} \frac{x - \sin x}{x + \sin x} = $ \blank{2cm}。

    \item 求极限 $ \lim_{x\to1}\left(\frac{x}{x-1}-\frac{1}{\ln x}\right)= $ \blank{2cm}。

    \item 求极限 $ \lim_{x\to\infty}x[\ln(x-2)-\ln(x+1)]= $ \blank{2cm}。

    \item 求极限 $ \lim_{x\to1}\frac{\arcsin(x^{2}-1)}{\ln x}= $ \blank{2cm}。

    \item 求极限 $ \lim_{x\to1}\frac{\sin(x^{3}-1)}{x-1}= $ \blank{2cm}。

    \item 求极限 $ \lim_{x\to0}\left(\frac{1}{2x}-\frac{1}{e^{2x}-1}\right) $。

    \item 求极限 $ \lim_{x \to 0} \left(\frac{1}{\sin^{2} x} - \frac{1}{x^{2}}\right) $。

    \item 求极限 $ \lim_{x \to e} \frac{x - e}{\ln x - 1} $。

    \item 求极限 $ \lim_{x\to0}\frac{\sqrt{1+\sin x}-\sqrt{1+\tan x}}{x^{3}} $。

    \item 求极限 $ \lim_{x\to0}\left(\lim_{n\to\infty}\cos\frac{x}{2}\cos\frac{x}{2^2}\cdots \cos\frac{x}{2^n}\right) $。
\end{enumerate}

\section*{第三章 连续（提高）}

\begin{enumerate}
    \item 已知函数 $ f(x) = \frac{x+2}{x^{2}-x} $ ，则 $ x = 0 $ 是函数 $ f(x) $ 的 \blank{1cm}。
    \choicesFour{可去间断点}{跳跃间断点}{无穷间断点}{连续点}

    \item 当 $ x \rightarrow +\infty $ 时， $ x\sin x $ 是 \blank{1cm}。
    \choicesFour{无穷大量}{无穷小量}{无界变量}{有界变量}

    \item 当 $ x \rightarrow 0 $ 时， $ x^{2}\sin\frac{1}{x} $ 是 \blank{1cm}。
    \choicesTwo
    {与 $x$ 等价无穷小}
    {$x$ 的低阶无穷小}
    {与 $x$ 等价无穷小}
    {$x$ 的高阶无穷小}

    \item 试确定当 $ x \rightarrow 0 $ 时，下列 \blank{1cm} 是关于 $x$ 的三阶无穷小。
    \choicesTwo
    {$ \sqrt[3]{x^{2}} - \sqrt{x} $}
    {$ \sqrt{a + x^{3}} - \sqrt{a} $ 其中 $ a \neq 0 $}
    {$ x^{3}+0.001x^{2} $}
    {$ \sqrt[3]{\tan x^{3}} $}

    \item 函数 $ f(x)=\frac{|x|\sin(x-1)}{x(x-1)(x-2)} $ 在下列区间有界的是 \blank{1cm}。
    \choicesFour{$ (0,1) $}{$ (1,2) $}{$ (0,2) $}{$ (2,3) $}

    \item 函数 $ y=\frac{1}{e^{\frac{x}{x-2}}-1} $ 的间断点为 \blank{2cm}。

    \item $ f(x)=\begin{cases}\sin\frac{1}{x}, & x>0\\ x-1, & x\leq0\end{cases} $ 函数 $ f(x) $ 的间断点是第 \blank{2cm} 类间断点。

    \item $ x = 1 $ 为函数 $ y = \frac{x^{4}-1}{x^{3}-1} $ 的 \blank{2cm} 间断点。

    \item 函数 $ f(x)=\frac{x^{3}-x}{\sin\pi x} $ 的可去间断点的个数为 \blank{2cm}。

    \item 若要使 $ f(x)=\frac{1-\sqrt{1-x}}{1-\sqrt[3]{1-x}} $ 在点 $ x=0 $ 处连续，则要补充定义 $ f(0)= $ \blank{2cm}。

    \item 设函数 $ f(x)=\begin{cases}x+a\cos2x, & x>0\\ e^{ax}-a, & x\leq0\end{cases} $ 为 $ (-\infty,+\infty) $ 上的连续函数，则 $ a= $ \blank{2cm}。

    \item 求极限 $ \lim_{x\to+\infty}\left(x+\sqrt{1+x^{2}}\right)^{\frac{1}{x}}= $ \blank{2cm}。

    \item 设函数 $ f(x)=\begin{cases}\frac{\sin2x+e^{2ax}-1}{x}, & x\neq0,\\ a, & x=0\end{cases} $ 在 $ (-\infty,+\infty) $ 上连续，则 $ a= $ \blank{2cm}。

    \item 证明方程 $ x^{5}+x-1=0 $ 只有一个正根。

    \item 设函数 $ f(x)=\begin{cases}x^{2}+1, & |x|\leq C\\ \frac{10}{|x|}, & |x|>C\end{cases} $ 在定义域上连续，试求常数 $C$。
\end{enumerate}

\section*{第四章 导数与微分（提高）}

\begin{enumerate}
    \item 设 $ f(0)=0 $ 且 $ \lim_{x\to0}\frac{f(x)}{x} $ 存在，则 $ \lim_{x\to0}\frac{f(x)}{x}= $ \blank{1cm}。
    \choicesFour{$ f'(x) $}{$ f'(0) $}{$ f(0) $}{$ \frac{1}{2}f'\left(\frac{1}{2}\right) $}

    \item 函数 $ y=f(x) $ 在点 $ x_{0} $ 处的左导数 $ f'_{-}(x_{0}) $ 和右导数 $ f'_{+}(x_{0}) $ 都存在,是 $ f(x) $ 在 $ x_{0} $ 可导的 \blank{1cm}。
    \choicesTwo
    {充分必要条件}
    {充分但非必要条件}
    {必要但非充分条件}
    {既非充分又非必要条件}

    \item 函数 $ f(x)=|\sin x| $ 在 $ x=0 $ 处 \blank{1cm}。
    \choicesTwo
    {可导}
    {连续但不可导}
    {不连续}
    {极限不存在}

    \item 设 $ y = \tan^{2}x $ , 则 $ dy = $ \blank{1cm}。
    \choicesTwo
    {$ 2\tan x \sec^{2} x dx $}
    {$ 2\sin x \cos^{2} x dx $}
    {$ 2\sec x \tan^{2} dx $}
    {$ 2\cos x \sin^{2} x dx $}

    \item 设 $ F(x)=f(x)+f(-x) $ ，且 $ f'(x) $ 存在，则 $ F'(x) $ 是 \blank{1cm}。
    \choicesTwo
    {奇函数}
    {偶函数}
    {非奇非偶的函数}
    {无法判定其奇偶性}

    \item 设周期函数 $ f(x) $ 在 $ (-\infty, \infty) $ 周期为 3, 则曲线 $ y = f(x) $ 在点 $ (4, f(4)) $ 的切线斜率为 \blank{2cm}。

    \item 若 $ y = e^{x}(\sin x + \cos x) $ ，则 $ \frac{dy}{dx} = $ \blank{2cm}。

    \item 曲线 $ xy + \ln y - 1 = 0 $ 在点 $(1,1)$ 处的法线方程是 \blank{2cm}。

    \item 已知 $ \begin{cases} x = a(\sin t - t\cos t) \\ y = a(\cos t + t\sin t) \end{cases} $ ，则 $ \left.\frac{dy}{dx}\right|_{t=\frac{3}{4}\pi} = $ \blank{2cm}。

    \item 若 $ f(x) = x^{2} \ln x $ , 则 $ f''(2) = $ \blank{2cm}。

    \item 对方程 $ e^{y} + xy = e^{2} $ 所确定的函数 $ y = y(x) $ 求 $ \frac{d^{2}y}{dx^{2}} $。

    \item 设函数 $ f(x)=\begin{cases}ax+b, & x<1\\ x^{2}, & x\geq1\end{cases} $ 在 $x=1$ 处可导，求 $a,b$ 的值。

    \item 求函数 $ y = \ln(x + \sqrt{x^{2} + a^{2}}) $ 的导数。

    \item 对方程 $ \ln(\sqrt{x^{2}+y^{2}})=\arctan\frac{y}{x} $ 所确定的函数 $ y=y(x) $ 求 $ \frac{d^{2}y}{dx^{2}} $。

    \item 求函数 $ y = \frac{x}{2} \sqrt{a^{2} - x^{2}} + \frac{a^{2}}{2} \arcsin \frac{x}{a} \quad (a > 0) $ 的导数。
\end{enumerate}

\section*{第五章 中值定理与导数（提高）}

\begin{enumerate}
    \item 曲线 $ y=\frac{1}{f(x)} $ 有水平渐近线的充分条件是 \blank{1cm}。
    \choicesTwo
    {$ \lim_{x\to\infty}f(x)=0 $}
    {$ \lim_{x\to\infty}f(x)=\infty $}
    {$ \lim_{x\to0}f(x)=0 $}
    {$ \lim_{x\to0}f(x)=\infty $}

    \item 曲线 $ y = \frac{x}{1 - x^{2}} $ 的渐近线有 \blank{1cm}。
    \choicesFour{1 条}{2 条}{3 条}{4 条}

    \item 点 $(1,2)$ 是 $ f(x)=(x-a)^{3}+b $ 对应图形的拐点，则 \blank{1cm}。
    \choicesFour{$ a=0,b=1 $}{$ a=2,b=3 $}{$ a=1,b=2 $}{$ a=-1,b=-6 $}

    \item 函数 $ y = |\ln x| $ 对应图形的拐点是 \blank{1cm}。
    \choicesFour{$ (1,0) $}{$ (e,1) $}{$ (2,\ln 2) $}{不存在}

    \item 设函数 $ f(x)=(x-1)(x-2)(x-3) $ ，则方程 $ f'(x)=0 $ 有 \blank{1cm}。
    \choicesFour{一个实根}{二个实根}{三个实根}{无实根}

    \item 设函数 $ f(x) = x^{3} + ax^{2} + bx + c $ , 且 $ f(0) = f'(0) = 0 $ , 则下列结论不正确的是 \blank{1cm}。
    \choicesTwo
    {$ b = c = 0 $}
    {当 $a > 0$ 时， $ f(0) $ 为极小值}
    {当 $a < 0$ 时， $ f(0) $ 为极大值}
    {当 $ a \neq 0 $ 时， $ (0, f(0)) $ 为拐点}

    \item 函数 $ y=\frac{x}{\ln x} $ 的单调增加的区间是 \blank{2cm}。

    \item 函数 $ y=(x-2)^{3} $ 的驻点是 \blank{2cm}。

    \item 函数 $ y=-x^{3}+x^{2} $ 的凸区间是 \blank{2cm}。

    \item 曲线 $ y = \frac{x^{2}-2x+2}{x-1} $ 的垂直渐近线的方程是 \blank{2cm}。

    \item 求曲线 $ y = x^{3} - 3x^{2} + 3 $ 的拐点 \blank{2cm}。

    \item 若 $ f(x) $ 在 $ [0,1] $ 上连续，在 $ (0,1) $ 内可导，且 $ f(0)=f(1)=0, f(1)=1 $ 。
    证明在 $ (0,1) $ 内至少有一点 $ \xi $ 使 $ f'(\xi)=1 $ 。

    \item 求斜边长为定长 $L$ 的直角三角形的最大面积。

    \item 讨论 $ y = xe^{-x} $ 的增减性、凹凸性、极值、拐点。

    \item 已知某产品的需求函数 $ Q = 28 - 2p $ 和总成本函数 $ C(Q) = 30 + 2Q $ , 其中 $Q$ 为销售量，$p$ 为价格，求当 $p$ 为多少可获得最大利润？

    \item 证明不等式 $ e^{x} < (1 + x)^{(1 + x)} $ ($x > 0$)。
\end{enumerate}

\section*{第六章 不定积分（提高）}

\begin{enumerate}
    \item 设函数 $ f(x) $ 的一个原函数为 $ \frac{1}{x} $ ，则 $ f(x)= $ \blank{1cm}。
    \choicesFour{$ -\frac{1}{x^{2}} $}{$ \frac{2}{x^{3}} $}{$ \frac{1}{x} $}{$ \ln|x| $}

    \item 若 $ e^{-x} $ 是 $ f(x) $ 的一个原函数，则 $ \int xf(x)dx = $ \blank{1cm}。
    \choicesTwo
    {$ e^{-x}(1-x)+C $}
    {$ e^{-x}(x+1)+C $}
    {$ e^{-x}(x-1)+C $}
    {$ -e^{-x}(x+1)+C $}

    \item 设 $ \int f(x)dx = \ln(x + \sqrt{1 + x^{2}}) + C, $ 则 $ f(x) = $ \blank{1cm}。
    \choicesFour{$ \frac{x}{\sqrt{1+x^{2}}} $}{$ \frac{1}{\sqrt{1+x^{2}}} $}{$ \frac{1}{1+x} $}{$ \frac{1}{1+x^{2}} $}

    \item $ \int 2^{3x} dx = $ \blank{1cm}。
    \choicesFour{$ \frac{2^{3x}}{3\ln 2} + C $}{$ \frac{1}{3}(\ln 2)2^{3x} + C $}{$ \frac{2^{x}}{3} + C $}{$ \frac{2^{3x}}{\ln 2} + C $}

    \item 设函数 $ f(x) $ 在 $ [a, b] $ 上连续， $ F(x) = \int_{0}^{x^{3}} f(t) dt $ ，则 $ F'(x) = $ \blank{1cm}。
    \choicesFour{$ f(x) $}{$ f(x^{3}) $}{$ 3x^{2}f(x^{3}) $}{$ 3x^{2}f(x) $}

    \item 设 $ f'(x^{2}) = \frac{1}{x} (x > 0) $ ，则 $ f(x) = $ \blank{2cm}。

    \item 设 $ f(x) $ 的一个原函数是 $ \sin x, $ 则 $ \int x f'(x)dx = $ \blank{2cm}。

    \item $ \int 3^{x}e^{x}dx = $ \blank{2cm}。

    \item 不定积分 $ \int\frac{x-3}{x^{2}+1}dx= $ \blank{2cm}。

    \item 不定积分 $ \int \frac{1}{e^{x} + e^{-x}} dx = $ \blank{2cm}。

    \item 已知 $ f(x) $ 的一个原函数为 $ \frac{\sin x}{1 + x \sin x} $ ，求 $ \int f(x) f'(x) dx $。

    \item 若 $ \int xf(x)dx = \arcsin x + C $ 求 $ I = \int \frac{1}{f(x)}dx $。

    \item 求不定积分 $ \int\frac{e^{\sqrt{x}}\sin\sqrt{x}}{2\sqrt{x}}dx $。

    \item 计算 $ \int \sin^{2} x \cos^{3} x dx $。

    \item 计算 $ \int \frac{\sqrt{\arcsin x}}{\sqrt{1-x^{2}}} dx $。
\end{enumerate}

\section*{第七章 定积分（提高）}

\begin{enumerate}
    \item $ \frac{d}{dx}[x\int_{0}^{x}\sqrt{1+t^{4}}dt]= $ \blank{1cm}。
    \choicesTwo
    {$ \int_{0}^{x}\sqrt{1+t^{4}}dt $}
    {$ 4x^{4}\int_{0}^{x}\sqrt{1+t^{4}}dt $}
    {$ \int_{0}^{x}\sqrt{1+t^{4}}dt + x\sqrt{1+x^{4}} $}
    {$ x\sqrt{1+x^{4}} $}

    \item 设 $ f(x) $ 是连续函数，则 $ \frac{d}{dx}\int_{x^{2}}^{-1}f(t)dt= $ \blank{1cm}。
    \choicesFour{$ f(x^{2}) $}{$ 2xf(x^{2}) $}{$ -f(x^{2}) $}{$ -2xf(x^{2}) $}

    \item $ f(x) $ 在 $ [a, b] $ 上连续是 $ \int_{a}^{b} f(x) \, dx $ 存在的 \blank{1cm}。
    \choicesTwo
    {充分必要条件}
    {充分条件非必要条件}
    {必要条件非充分条件}
    {既非必要也非充分条件}

    \item $ \frac{d}{dx}\int_{0}^{x}g(x)f(t)dt= $ \blank{1cm}。
    \choicesTwo
    {$ g(x)f(x) $}
    {$ g'(x)f'(x) $}
    {$ g'(x)f(x) + g(x)f'(x) $}
    {$ g(x)f(x) + g'(x) \int_{0}^{x} f(t) dt $}

    \item 设函数 $ y = \int_{0}^{x} (t - 1) dt $ ，则 $y$ 有 \blank{1cm}。
    \choicesFour{极小值 $ \frac{1}{2} $}{极小值 $ -\frac{1}{2} $}{极大值 $ \frac{1}{2} $}{极大值 $ -\frac{1}{2} $}

    \item 已知 $ f(x) $ 连续，则 $ \frac{d}{dx} \int_{0}^{x} f(x - t) dt = $ \blank{2cm}。

    \item $ \int_{1}^{e^{3}}\frac{dx}{x\sqrt{1+\ln x}}= $ \blank{2cm}。

    \item 定积分 $ \int_{0}^{2}\sqrt{4-x^{2}}dx= $ \blank{2cm}。

    \item 定积分 $ \int_{-1}^{1}\frac{x^{2}+x^{5}\sin x^{2}}{1+x^{2}}dx= $ \blank{2cm}。

    \item 定积分 $ \int_{0}^{a}x^{2}\sqrt{a^{2}-x^{2}}dx= $ \blank{2cm}。

    \item 设函数 $ f(x)=\lim_{n\to\infty}\frac{d}{dx}\left(\int_{-1}^{x}\frac{1+t}{1+t^{4n}}dt\right) $ ，求 $ f(x) $ 的间断点。

    \item 求定积分 $ \int_{\frac{1}{e}}^{e}|\ln x|dx $。

    \item 已知 $ f(x) = e^{x^{2}} $ ，求 $ \int_{0}^{1} f'(x) f''(x) dx $。

    \item 求定积分 $ \int_{0}^{\frac{1}{2}}\arcsin x dx $。

    \item 设 $ f(x) $ 在 $ [a,b] $ 上连续，且单调增加，证明：
    $$ \int_{a}^{b}tf(t)dt\geq\frac{a+b}{2}\int_{a}^{b}f(t)dt $$
\end{enumerate}

\section*{第八章 微分方程（提高）}

\begin{enumerate}
    \item 微分方程 $ x \ln x dy + (y - \ln x) dx = 0 $ 满足 $ y|_{x=e} = 1 $ 的特解为 \blank{1cm}。
    \choicesTwo
    {$ \frac{1}{2} (\ln x + \frac{1}{\ln x}) $}
    {$ \frac{1}{2} (x + \frac{1}{\ln x}) $}
    {$ \frac{1}{2} (\ln x + \frac{1}{x}) $}
    {$ \frac{1}{2} (x + \frac{1}{x}) $}

    \item 微分方程 $ xy' + y = \frac{1}{1 + x^{2}} $ 满足 $ y|_{x=\sqrt{3}} = \frac{\sqrt{3}}{9} $ 的解在 $ x = 1 $ 处的值为 \blank{1cm}。
    \choicesFour{$ \frac{\pi}{4} $}{$ \frac{\pi}{3} $}{$ \frac{\pi}{2} $}{$ \pi $}

    \item 曲线 $ y = x^{2} $ 与 $ y^{2} = x $ 围成的平面图形绕 $y$ 轴旋转一周生成的立体体积为 \blank{1cm}。
    \choicesFour{$ \frac{1}{3} $}{$ \frac{\pi}{2} $}{$ \frac{1}{2} $}{$ \frac{3\pi}{10} $}

    \item 若 $ C_{1} $ 和 $ C_{2} $ 为两个独立的任意常数，则 $ y = C_{1} \cos x + C_{2} \sin x $ 为下列哪个方程的通解 \blank{1cm}。
    \choicesFour{$ y''+y=0 $}{$ y''+y=x^{2} $}{$ y''-3y'+2y=0 $}{$ y''+y'-2y=2x $}

    \item 微分方程 $ (y'')^{5}+2y'+xy^{6}=0 $ 的阶数是 \blank{1cm}。
    \choicesFour{1}{2}{3}{4}

    \item 已知函数 $ y = y(x) $ 在任意点处的增量 $ \Delta y = \frac{y \Delta x}{1 + x^{2}} + \alpha $, 且当 $ \Delta x \to 0 $ 时， $ \alpha $ 为 $ \Delta x $ 的高阶无穷小，若 $ y(0) = \pi $ , 则 $ y(1) = $ \blank{2cm}。

    \item 求由曲线 $ y = x^{2} $ 及 $ y = \sqrt{x} $ 所围成的平面图形的面积 \blank{2cm}。

    \item 已知 $ y = e^{x}(C_{1} \cos \sqrt{2} x + C_{2} \sin \sqrt{2} x) $ ($C_{1}, C_{2}$ 为任意常数）是某二阶常系数线性微分方程的通解，其对应的方程为 \blank{2cm}。

    \item 求微分方程 $ xy' + y = xe^{x} $ 满足 $ y|_{x=1} = 1 $ 的特解 \blank{2cm}。

    \item 求微分方程 $ \frac{dy}{dx}-\frac{y}{x}=x\sin x $ 的通解 \blank{2cm}。

    \item 求微分方程 $ y'' + 3y' = 3x $ 的通解。

    \item 求微分方程 $ y''-2y'+y=2e^{x} $ 的通解。

    \item 求微分方程 $ y'+2xy=xe^{-x^{2}} $ 满足初始条件 $ y|_{x=0}=1 $ 的特解。

    \item 求微分方程 $ y' - y \cot x = 2x \sin x $ 的通解。

    \item 计算抛物线 $ y^{2} = 2x $ 与直线 $ y = x - 4 $ 所围成的图像的面积。
\end{enumerate}

\section*{第九章 多元微积分（提高）}

\begin{enumerate}
    \item 已知函数 $ z=\frac{\sin(xy)}{y} $ ，则 $ \frac{\partial^{2}z}{\partial x\partial y}= $ \blank{1cm}。
    \choicesFour{$ -x\sin(xy) $}{$ x\sin(xy) $}{$ -x\cos(xy) $}{$ x\cos(xy) $}

    \item 二元函数 $ f(x,y) $ 在点 $ (x_{0},y_{0}) $ 处存在偏导数是 $ f(x,y) $ 在该点可微的 \blank{1cm}。
    \choicesTwo
    {必要而不充分条件}
    {充分而不必要条件}
    {必要且充分条件}
    {既不必要也不充分条件}

    \item 函数 $ f(x,y) $ 可微， $ f'_{x}(x_{0},y_{0})=0, f'_{y}(x_{0},y_{0})=0 $ ，则函数 $ f(x,y) $ 在 $ (x_{0},y_{0}) $ \blank{1cm}。
    \choicesTwo
    {可能有极值，也可能没有极值}
    {必有极值，可能是极大值也可能是极小值}
    {一定没有极值}
    {必有极值，且为极小值}

    \item $ \lim_{x\to0}\frac{xy^{2}}{x^{2}+y^{4}}= $ \blank{1cm}。
    \choicesFour{0}{1}{$ \frac{1}{2} $}{不存在}

    \item $ f_{x}^{'}(x,y) $ 和 $ f_{y}^{'}(x,y) $ 在点 $ (x_{0},y_{0}) $ 连续是 $ f(x,y) $ 在点 $ (x_{0},y_{0}) $ 可微分的 \blank{1cm}。
    \choicesFour{充分条件}{必要条件}{充要条件}{无关条件}

    \item 设 $ z = \arctan \frac{y}{x} $ ，则 $ \frac{\partial z}{\partial x} = $ \blank{2cm}。

    \item 设 $ z = \ln\sqrt{x^{2} + y^{2}} $ , 则 $ x\frac{\partial z}{\partial x} + y\frac{\partial z}{\partial y} = $ \blank{2cm}。

    \item 设 $ z = \ln(\sqrt{x} + \sqrt{y}) $ ，则 $ x \frac{\partial z}{\partial x} + y \frac{\partial z}{\partial y} = $ \blank{2cm}。

    \item 二元函数 $ f(x,y)=x^{2}(2+y^{2})+y\ln y $ 的极小值 \blank{2cm}。

    \item $ z = \ln(x^{2} + y^{2} - 2) + \sqrt{4 - x^{2} - y^{2}} $ 函数的定义域为 \blank{2cm}。

    \item 已知函数 $ z = f(u, v) $ 可导， $ u = x \arcsin y, v = \frac{y}{x} $ ，求 $ \frac{\partial z}{\partial x}, \frac{\partial z}{\partial y} $。

    \item 求内接于半径为 $R$ 的球体且有最大体积的长方体的边长。

    \item 设函数 $ z = z(x, y) $ 由方程 $ z^{3} - 3xyz = a^{3} $ 确定，求 $ \frac{\partial z}{\partial x}, \frac{\partial z}{\partial y} $ 和 $dz$。

    \item 设 $ z = f(2x - y) + g(x, xy) $ ，其中函数 $ f(w) $ 具有二阶导数， $ g(u, v) $ 具有二阶连续偏导数，求 $ \frac{\partial z}{\partial x} $ 与 $ \frac{\partial^{2}z}{\partial x \partial y} $。

    \item 求内接于椭圆 $ \frac{x^{2}}{3}+\frac{y^{2}}{4}=1 $ 的最大长方形的面积。
\end{enumerate}

\section*{第十章 向量代数与空间解析几何（提高）}

\begin{enumerate}
    \item 直线 $ L_{1} $ : $ x - 1 = \frac{y - 5}{-2} = z + 8 $ 与直线 $ L_{2} $ : $ \begin{cases} x - y = 6 \\ 2y + z = 3 \end{cases} $ 的夹角为 \blank{1cm}。
    \choicesFour{$ \frac{\pi}{6} $}{$ \frac{\pi}{4} $}{$ \frac{\pi}{3} $}{$ \frac{\pi}{2} $}

    \item 曲面 $ z = 4 - x^{2} - y^{2} $ 上一点 $P$ 的切平面平行于平面 $ 2x + 2y + z - 1 = 0 $ ，则 $P$ 点坐标为 \blank{1cm}。
    \choicesFour{$ (1,-1,2) $}{$ (-1,1,2) $}{$ (1,1,2) $}{$ (-1,-1,2) $}

    \item 直线 $l$: $ \frac{x+3}{-2}=\frac{y+4}{-7}=\frac{z}{3} $ 与平面 $ \pi $: $4x-2y-2z-3=0$ 的位置关系是 \blank{1cm}。
    \choicesFour{平行}{垂直相交}{$l$ 在 $ \pi $ 上}{相交但不垂直}

    \item 已知向量 $\mathbf{a}$ 垂直于向量 $ \mathbf{b} = 2\mathbf{i} - 3\mathbf{j} + \mathbf{k} $ 和 $ \mathbf{c} = \mathbf{i} - 2\mathbf{j} + 3\mathbf{k} $ ，且满足于 $ \mathbf{a} \cdot (\mathbf{i} + 2\mathbf{j} - 7\mathbf{k}) = 10 $ ，求 $\mathbf{a} =$ \blank{1cm}。
    \choicesTwo
    {$ -7\mathbf{i} - 5\mathbf{j} - \mathbf{k} $}
    {$ 7\mathbf{i} + 5\mathbf{j} + \mathbf{k} $}
    {$ -5\mathbf{i} - 3\mathbf{j} - \mathbf{k} $}
    {$ 5\mathbf{i} + 3\mathbf{j} + \mathbf{k} $}

    \item 向量 $\mathbf{a}$ 与 $x$ 轴与 $y$ 轴构成等角，与 $z$ 轴夹角是前者的 2 倍，下面哪一个代表的是 $\mathbf{a}$ 的方向 \blank{1cm}。
    \choicesTwo
    {$ \alpha=\frac{\pi}{2},\beta=\frac{\pi}{2},\gamma=\frac{\pi}{4} $}
    {$ \alpha=\frac{\pi}{4},\beta=\frac{\pi}{4},\gamma=\frac{\pi}{8} $}
    {$ \alpha=\frac{\pi}{4},\beta=\frac{\pi}{4},\gamma=\frac{\pi}{2} $}
    {$ \alpha=\pi,\beta=\frac{\pi}{2},\gamma=\frac{\pi}{2} $}

    \item 设已知两点 $A(4,0,5)$ 和 $B(7,1,3)$，方向和 $ \overrightarrow{AB} $ 一致的单位向量为 \blank{2cm}。

    \item 已知 $ \mathbf{a} = \{2, 1, 2\} $ , $ \mathbf{b} = \{4, -1, 10\} $ , $ \mathbf{c} = \mathbf{b} - \lambda \mathbf{a} $ , 且 $ \mathbf{a} \perp \mathbf{c} $ , 则 $ \lambda = $ \blank{2cm}。

    \item 过点 $ (3,2,1) $ 与直线 $ x = y = z $ 垂直的平面方程为 \blank{2cm}。

    \item 过点 $ (3,-2,2) $ 与直线 $ x=\frac{y}{2}=z $ 垂直的平面方程是 \blank{2cm}。

    \item 过点 $ (2,2,3) $ 且与平面 $ 3x + y = z + 2 $ 垂直的直线方程是 \blank{2cm}。

    \item 直线 $ \frac{x-1}{2} = \frac{y+2}{5} = \frac{z-4}{-3} $ 的方向向量 $s=$ \blank{1.5cm}，与平面 $ 2x + 5y - 3z - 4 = 0 $ 是 \blank{1.5cm} 的。

    \item 求平行于 $y$ 轴且过点 $ P(1,5,1) $ , $ Q(3,2,-1) $ 的平面方程。

    \item 已知 $ \triangle ABC $ 的三个顶点的坐标分别为 $ A(1,2,3) $ , $ B(3,4,5) $ , $ C(2,4,7) $ ,
    求：(1) $ \triangle ABC $ 所在的平面方程；
    (2) $AB$ 边上的中线 $CD$ 所在的直线方程。

    \item 求过点 $ (-1,-4,3) $ 并与两直线 $ L_{1}\begin{cases}2x-4y+z=1\\ x+3y+5=0\end{cases} $ 和 $ L_{2}\begin{cases}x=2+4t\\ y=-1-t\\ z=-3+2t\end{cases} $ 都垂直的直线方程。

    \item 求通过点 $ M_{1}(3,-5,1) $ ， $ M_{2}(4,1,2) $ 且垂直于平面 $ x-8y+3z-1=0 $ 的平面方程。
\end{enumerate}

\section*{第十一章 二重积分（提高）}

\begin{enumerate}
    \item 设函数 $ f(x,y) $ 连续，则 $ \int_{1}^{2}dy\int_{y}^{2}f(x,y)dx $ 交换积分次序为 \blank{1cm}。
    \choicesTwo
    {$ \int_{1}^{2}dx\int_{1}^{x}f(x,y)dy $}
    {$ \int_{1}^{2}dx\int_{x}^{1}f(x,y)dy $}
    {$ \int_{1}^{2}dx\int_{2}^{x}f(x,y)dy $}
    {$ \int_{1}^{2}dx\int_{x}^{2}f(x,y)dy $}

    \item 设 $D$ 是方形域： $ 0 \leq x \leq 1, 0 \leq y \leq 1, \iint_{D} xyd\sigma = $ \blank{1cm}。
    \choicesFour{1}{$ \frac{1}{2} $}{$ \frac{1}{3} $}{$ \frac{1}{4} $}

    \item $ I = \int_{0}^{1} dy \int_{-\sqrt{1-y}}^{\sqrt{1-y}} 3x^{2}y^{2} dx $ 交换积分次序为 $I =$ \blank{1cm}。
    \choicesTwo
    {$ \int_{0}^{1}dy\int_{0}^{\sqrt{1-y}}3x^{2}y^{2}dx $}
    {$ \int_{0}^{\sqrt{1-y}}dy\int_{0}^{1}3x^{2}y^{2}dx $}
    {$ \int_{-1}^{1}dx\int_{0}^{1-x^{2}}3x^{2}y^{2}dy $}
    {$ \int_{0}^{1}dx\int_{0}^{1+x^{2}}3x^{2}y^{2}dy $}

    \item 二重积分 $ \iint_{D} f(x,y)d\sigma $ 在极坐标下的面积元素为 \blank{1cm}。
    \choicesFour{$ d\sigma = dx dy $}{$ d\sigma = r dr d\theta $}{$ d\sigma = dr d\theta $}{$ d\sigma = r^{2}\sin\theta dr d\theta $}

    \item [5.] (题目序号更正) 设 $ I = \int_{0}^{1} dx \int_{x^{2}}^{x} f(x, y) dy $ 交换积分顺序后，则 $ I = $ \blank{2cm}。

    \item [6.] $D$ 是圆环 $ a^{2}\leq x^{2}+y^{2}\leq b^{2} $ ，计算 $ I=\iint_{D}(x^{2}+y^{2})d\sigma $。

    \item [7.] 交换积分顺序： $ \int_{0}^{\frac{1}{4}}dy\int_{y}^{\sqrt{y}}f(x,y)dx+\int_{\frac{1}{4}}^{\frac{1}{2}}dy\int_{y}^{\frac{1}{2}}f(x,y)dx $ 为 \blank{2cm}。

    \item [8.] 计算二重积分 $ \iint_{D} \frac{y}{x} d\sigma $ ，其中 $D$ 由 $ x^{2} + y^{2} \leq a^{2} (a > 0) $ ，$y = x$ 及 $x$ 轴在第一象限所围成的区域。

    \item [9.] 二重积分 $ I=\iint_{D}\ln(1+x^{2}+y^{2})dxdy $ ,其中 $D$ 是由圆周 $ x^{2}+y^{2}=1 $ 及坐标轴所围成的第一象限内的区域。

    \item [10.] 求二重积分 $ \iint_{D} (x + 2y) \, dx \, dy $ ，其中区域 $D$ 是由 $ y = x, y = 4x, xy = 1 $ 围成的第一象限部分的闭区域。

    \item [11.] 求 $ \iint_{D} ye^{xy} dx dy $ ，其中 $D$ 是由 $ xy = 1, x = 2, y = 1 $ 所围成的区域。

    \item [12.] 求二重积分 $ \iint_{D}\frac{y}{x}dxdy $ ，其中 $D$ 是由 $ x=e,y=\ln x,y=0 $ 所围成的区域。

    \item [13.] 计算 $ \iint_{D} \frac{x}{y} d\sigma $ ，其中 $D$ 是由 $ y = 1, y = x^{2}, x = 2 $ 所围成的闭区域。

    \item [14.] 计算二重积分 $ \iint_{D} x\sqrt{y} $ 其中 $D$ 是由 $ y = \sqrt{x}, y = x^{2} $ 围成的区域。
\end{enumerate}

\section*{第十二章 无穷级数提高十五题}

\begin{enumerate}
    \item 有下列关于数项级数的命题，其中正确的个数为 \blank{1cm}。
    \begin{enumerate}[label=(\arabic*)]
        \item 若 $ \lim_{n\to\infty}u_{n}\neq0 $ ,则 $ \sum_{n=1}^{\infty}u_{n} $ 必发散.
        \item 若 $ u_{n}\geq0,u_{n}\geq u_{n+1}(n=1,2,3\ldots) $ 且 $ \lim_{n\to\infty}u_{n}=0, $ 则 $ u_{n} $ 必收敛.
        \item 若 $ \sum_{n=1}^{\infty}u_{n} $ 收敛，则 $ \sum_{n=1}^{\infty}|u_{n}| $ 必收敛.
        \item 若 $ \sum_{n=1}^{\infty}u_{n} $ 收敛于 $s$,则任意改变该级数项的位置所得到的新的级数仍收敛于 $s$.
    \end{enumerate}
    \choicesFour{0}{1}{2}{3}

    \item 下列级数中收敛的级数是 \blank{1cm}。
    \choicesFour{$ \sum_{n=1}^{\infty}\frac{1}{3^{n}} $}{$ \sum_{n=1}^{\infty}\frac{1}{n+1} $}{$ \sum_{n=1}^{\infty}\frac{3^{n}}{2^{n}} $}{$ \sum_{n=1}^{\infty}\frac{1}{\sqrt{n+1}} $}

    \item 级数 $ \sum_{n=1}^{\infty} (-1)^{n} \frac{1}{n^{2}} $ 是 \blank{1cm}。
    \choicesFour{发散的}{条件收敛}{绝对收敛}{收敛性不能确定}

    \item 下列级数中为条件收敛的级数是 \blank{1cm}。
    \choicesTwo
    {$ \sum_{n=1}^{\infty}(-1)^{n}\frac{n}{n+1} $}
    {$ \sum_{n=1}^{\infty}(-1)^{n}\sqrt{n} $}
    {$ \sum_{n=1}^{\infty}(-1)^{n}\frac{1}{n^{2}} $}
    {$ \sum_{n=1}^{\infty}(-1)^{n}\frac{1}{\sqrt{n}} $}

    \item 幂级数 $ \sum_{n=0}^{\infty}\frac{x^{n}}{n!} $ 的收敛半径是 \blank{1cm}。
    \choicesFour{1}{0}{$ +\infty $}{$ \frac{1}{2} $}

    \item 求幂级数 $ \sum_{n=1}^{\infty} \frac{(x-5)^{n}}{\sqrt{n}} $ 的收敛域 \blank{2cm}。

    \item 判别级数的收敛性，级数是 $ \sum_{n=1}^{\infty}\frac{n^{n}}{(n!)^{2}} $ \blank{2cm}。(收敛还是发散)

    \item 判别级数的收敛性，级数是 $ \sum_{n=0}^{\infty} (-1)^{n} \frac{(n+1)!}{n^{n+1}} $ \blank{2cm}。(绝对收敛，条件收敛还是发散)

    \item 求幂级数 $ \sum_{n=0}^{\infty} 10^{n}(x-1)^{n} $ 的收敛域 \blank{2cm}。

    \item 求幂级数 $ \sum_{n=0}^{\infty} \frac{x^{n}}{2^{n-1}} $ 的和函数 \blank{2cm}。

    \item 求幂级数 $ \sum_{n=0}^{\infty}\frac{x^{n}}{\sqrt{n+1}} $ 的收敛域。

    \item 求幂级数 $ \sum_{n=0}^{\infty} \frac{x^{n+2}}{n+1} $ 的收敛域及和函数。

    \item 求幂级数 $ \sum_{n=1}^{\infty} n(x-1)^{n} $ 的收敛域及和函数。

    \item 求幂级数 $ \sum_{n=2}^{\infty}\frac{1}{(n^{2}-1)2^{n}} $ 的收敛域及和函数。

    \item 求幂级数 $ \sum_{n=1}^{\infty} (-1)^{n} \frac{n(n+1)}{2^{2n}} $ 的收敛域及和函数。
\end{enumerate}

\end{document}